\chapter{Why Zeno Is as It Is}\label{ch-design}

Every chapter of this book has had boxes that give the reason for a
choice. This last chapter steps back and looks at the choices
together: the principles they come from, what Zeno took from the
languages before it, what it left out on purpose, and what it does
not have yet.

\section{Unambiguous and intuitive}

One sentence was set down at the start, and every feature was held to
it: \emph{the entire language must be unambiguous and intuitive}.
Unambiguous means a line has one reading, and a reader never has to
work out which of two things a symbol means. Intuitive means the
reading it has is the one a newcomer would guess. When the two pulled
apart, a message from the compiler bridged them: \verb|elif|,
\verb|&&|, \verb|def| and \verb|;| are each met with the Zeno spelling.

\section{The five principles}

Most decisions fall out of five principles. When two choices are
otherwise equal, the one higher on the list wins.

\begin{enumerate}
\item \textbf{Catch the mistake as early as possible.} Before the
  program runs is better than while it runs; while it runs is better
  than a wrong answer. Zeno should never quietly do something the
  programmer did not mean.
\item \textbf{Read like plain words.} Code is read far more often than
  it is written. A whole word beats a symbol or an abbreviation.
\item \textbf{Punctuation has to earn its place.} A mark stays only if
  removing it would make code harder to read or easier to get wrong.
\item \textbf{One way to say a thing, with one meaning.} An operator or
  keyword should not change meaning with the values it is given.
\item \textbf{Nothing a microcontroller cannot afford.} Every feature
  must compile to small, predictable C with no hidden runtime.
\end{enumerate}

Where Python already had a spelling that passed these tests, Zeno used
it, because more people can read Python than any other language.

\section{What came from where}

Zeno set out to blend the good parts of the big languages. Here is
where each part came from:

\begin{center}
\begin{tabular}{lP{9.4cm}}
\toprule
From & Zeno took \\
\midrule
Python & \verb|and|/\verb|or|/\verb|not|, \verb|in|, \verb|try|/\verb|except|/\verb|raise|,
  f-strings, comprehensions, \verb|import| and \verb|from|, \verb|//| and \verb|**|,
  the precedence of the bit operators, chained comparisons \\
C & braces, \verb|0x| numbers, the bit operators, compiling to C itself,
  integer promotion (without its silent cut) \\
C++ & generics checked for each use, compiled to a copy for each type \\
Rust & \verb|let|, one owner for a value, \verb|unsafe| blocks, the sized
  number names, enums that carry data and a \verb|match| that must be complete \\
Java & \verb|interface|, and a class that says which ones it is \\
JavaScript & functions as values with closures, written in place \\
Go and Swift & no semicolons; Go's square brackets for generics and its one
  layout for every file \\
Kotlin and TypeScript & a checker that follows \verb|if x is not none| \\
\bottomrule
\end{tabular}
\end{center}

\section{What was left out, and why}

As much of Zeno's design is in what it does not have:

\begin{description}
\item[No null that can appear anywhere.] Only \verb|T or none| can be
  \verb|none|, and it must be asked about.
\item[No quiet conversions.] An \verb|int| does not become a
  \verb|float|, a \verb|bool| or a \verb|str| by itself, and nothing is
  true or false but a \verb|bool|.
\item[No wrapping arithmetic by default.] An overflow is an error;
  wrapping is a word, \verb|wrap|.
\item[No negative indexes.] A position below zero is a mistake, not the
  other end of the list.
\item[No inheritance.] A class is read in one place; interfaces do
  what inheritance was mostly for.
\item[No constructor of your own.] The fields are the constructor, and
  a named function builds what needs checking.
\item[No exceptions that unwind.] Errors are declared with
  \verb|raises| and compile to jumps or to return values, with no
  tables and no cost where they are not used.
\item[No mutable default arguments shared between calls,] no
  \verb|is| for comparing values, no \verb|elif|, no \verb|lambda|, no
  \verb|=>|.
\end{description}

Each of these is something a language before Zeno had, and that
somebody, somewhere, lost an afternoon to.

\section{How decisions were made}

Every decision in Zeno has a written reason, in the file
\verb|docs/decisions.md| of the repository: what was decided, why, and
what was turned down. A decision made on less than it deserved was
marked \emph{provisional}. Before version 0.1 each of the eighty or so
was read again and given one of three answers: settled, a stated limit
of 0.1, or changed. A choice with no good reason should not be in the
language, and if a better reason turns up for the opposite choice, the
choice should change.

The decisions are numbered, and the boxes in this book cite them, so
anyone who wants the full argument for a rule can find it.

\section{What Zeno does not have yet}

Every limit is refused today with a message, so each can be added later
without changing what any program written now means. The main ones:

\begin{itemize}
\item \textbf{On a board:} reading a pin, an analogue converter or a
  serial port, interrupts and timers. These must be designed against
  real hardware.
\item \textbf{Numbers:} a \verb|u64|, units made from other units, and
  reading a file as bytes.
\item \textbf{Types:} unions other than \verb|or none|, interfaces with
  type parameters, and bounds on generics.
\item \textbf{The library:} reading a file
  in pieces, \verb|program.exit|, the time of day.
\item \textbf{Builds:} compiling for another processor in one step,
  and a few corners such as a named function inside a function.
\item \textbf{The web:} images on a canvas, touch, and a page small
  and fast enough for a game.
\end{itemize}

\hyperref[app-limits]{Appendix G} has the full list.

\section{The name}

Zeno is named after Zeno of Elea, whose paradoxes say that Achilles can
never catch the tortoise: before he reaches it, he must reach halfway,
and before that, halfway to halfway, for ever. And after the quantum
Zeno effect, where a system watched closely enough never changes. Its
motto is \emph{every journey is halfway done, forever}: a fitting name
for a language that will always be almost finished, and a good thing
for a language to be, since it means there is always the next half to
build.

\begin{summary}
Zeno is held to one sentence, unambiguous and intuitive, and five
principles: catch mistakes early, read like words, earn every mark, one
way with one meaning, nothing a microcontroller cannot afford. It takes
from Python, C, C++, Rust, Java, JavaScript, Go, Swift, Kotlin and
TypeScript, leaves out what has cost programmers most, and writes the
reason for every choice down.
\end{summary}
